\documentclass[11pt]{article}

\usepackage[margin=1in]{geometry}
\usepackage{amsmath, amssymb}
\usepackage{xcolor}
\usepackage{listings}
\usepackage{hyperref}
\usepackage{array}
\usepackage{booktabs}

\hypersetup{
    colorlinks=true,
    linkcolor=blue!60!black,
    urlcolor=blue!60!black,
    citecolor=blue!60!black
}

\lstdefinestyle{bugpython}{
  language=Python,
  basicstyle=\ttfamily\small,
  keywordstyle=\color{blue!60!black},
  stringstyle=\color{green!40!black},
  commentstyle=\color{gray!70!black},
  showstringspaces=false,
  breaklines=true,
  frame=single,
  rulecolor=\color{gray!30},
  columns=fullflexible,
  keepspaces=true
}

\title{SymPy Bug Report}
\author{}
\date{}

\begin{document}
\setlength{\emergencystretch}{2em}

\maketitle

\section{Bug 38: Spurious Imaginary Part in a Divergent Real Integral}

\subsection{Minimal Reproducer}

\medskip

\begin{lstlisting}[style=bugpython]
from sympy import Abs, Rational, integrate, symbols

x = symbols("x", real=True)
val = integrate(1/Abs(x), (x, -1, 0))
print("integral =", val)
print("truncated eps=1/10:",
      integrate(1/Abs(x), (x, -1, -Rational(1, 10))))
print("truncated eps=1/100:",
      integrate(1/Abs(x), (x, -1, -Rational(1, 100))))
\end{lstlisting}

\subsection{Observed Behavior}

\medskip

\begin{lstlisting}[style=bugpython]
integral = oo + I*pi
truncated eps=1/10: log(10)
truncated eps=1/100: log(100)
\end{lstlisting}

SymPy reports the improper integral as \(\infty + i\pi\).  The finite
truncations are real and positive, and the integrand is real and positive on
the interval of integration except at the singular endpoint.  The imaginary
term is therefore not a harmless alternate representation of the real
improper integral.

\subsection{Expected Behavior}

The mathematically correct result is positive real infinity:
\[
  \int_{-1}^{0} \frac{1}{|x|}\,dx = \infty.
\]
The result should not contain any finite imaginary component.

\subsection{Mathematical Explanation}

For \(x \in [-1,0)\), \(|x|=-x\), so
\[
  \frac{1}{|x|} = -\frac{1}{x}.
\]
The corresponding truncated improper integral is
\[
  \int_{-1}^{-\varepsilon} \frac{1}{|x|}\,dx
  =
  \int_{-1}^{-\varepsilon} -\frac{1}{x}\,dx
  =
  -\log(\varepsilon),
  \qquad \varepsilon>0.
\]
As \(\varepsilon \to 0^+\), this tends to \(+\infty\).  More generally,
\[
  \int_{-a}^{-\varepsilon} \frac{1}{|x|}\,dx
  = \log\!\left(\frac{a}{\varepsilon}\right)
  \to +\infty
  \quad (a>0).
\]
The integral is over a real-valued, nonnegative function on the real line.
No branch choice for the complex logarithm should contribute a residual
\(i\pi\) to this real improper integral.

\subsection{Source-Level Diagnosis}

The diagnosis identifies the relevant public integration path as
\texttt{Integral.doit} in \nolinkurl{sympy/integrals/integrals.py}.  Because
the integrand contains \(\operatorname{Abs}\), the integrator rewrites
\(1/\operatorname{Abs}(x)\) as a real \texttt{Piecewise} expression, equivalent
to
\[
  \operatorname{Piecewise}\bigl((1/x, x\ge 0),(-1/x,\mathrm{True})\bigr).
\]
The \texttt{Piecewise} antiderivative path then produces an antiderivative
equivalent to
\[
  \operatorname{Piecewise}\bigl((-\log(x), x\le 0),(\log(x),\mathrm{True})\bigr).
\]
Definite interval evaluation proceeds through \texttt{Integral.doit} in
\nolinkurl{sympy/integrals/integrals.py}, which calls \texttt{\_eval\_interval}
on the antiderivative.  The diagnosis locates the responsible endpoint
behavior in \texttt{\_eval\_interval} in \nolinkurl{sympy/core/expr.py}, used
through \texttt{Piecewise.\_eval\_interval} in
\nolinkurl{sympy/functions/elementary/piecewise.py}.

The problematic rule is that endpoint evaluation first substitutes the
endpoint into the antiderivative and only falls back to a one-sided limit if
that substituted value is infinite or not a number:

\begin{lstlisting}[style=bugpython]
C = self.subs(x, c)
\end{lstlisting}

At the finite endpoint \(x=-1\), substituting into the negative-branch
antiderivative \(-\log(x)\) gives \(-i\pi\), using the principal complex
logarithm.  Since this is finite, \texttt{\_eval\_interval} keeps it.  At the
singular endpoint \(x \to 0^-\), the one-sided limit is \(\infty\).  The
interval evaluation therefore forms
\[
  \infty - (-i\pi) = \infty + i\pi.
\]

This source-level behavior evaluates a real-branch antiderivative using the
principal complex value of \(\log(x)\) at a negative real endpoint.  On the
negative real interval, an antiderivative for \(-1/x\) should behave like
\(-\log(-x)\), or equivalently use a real branch value up to a real constant.
A plausible fix direction supported by the diagnosis is to evaluate
\texttt{Piecewise} antiderivative endpoints by one-sided limits within the
active real branch, or to represent this negative-branch antiderivative in a
real-valued form such as \(-\log(-x)\) or \(-\log(|x|)\).

\subsection{Additional Failing Instances}

The same failure occurs for the family
\[
  \int_{-a}^{0} \frac{1}{|x|}\,dx
  \qquad (a \in \{1,2,3,4,5,6\}).
\]
For every positive \(a\), the integrand is real and nonnegative on
\([-a,0)\), and the truncated integral
\(\int_{-a}^{-\varepsilon} 1/|x|\,dx=\log(a/\varepsilon)\) diverges to
\(+\infty\).  Thus the expected result is \(\infty\) with no imaginary part
throughout the family.

\begin{center}
\small
\renewcommand{\arraystretch}{1.3}
\begin{tabular}{@{}>{\raggedright\arraybackslash}p{0.44\linewidth}>{\raggedright\arraybackslash}p{0.23\linewidth}>{\raggedright\arraybackslash}p{0.23\linewidth}@{}}
\toprule
\textbf{Input / Expression} & \textbf{SymPy behavior} & \textbf{Expected behavior} \\
\midrule
\(\int_{-1}^{0} 1/|x|\,dx\) & \(\infty+i\pi\) & \(\infty\) \\
\(\int_{-2}^{0} 1/|x|\,dx\) & not equal to \(\infty\); includes a spurious imaginary branch contribution & \(\infty\) \\
\(\int_{-3}^{0} 1/|x|\,dx\) & not equal to \(\infty\); includes a spurious imaginary branch contribution & \(\infty\) \\
\(\int_{-4}^{0} 1/|x|\,dx\) & not equal to \(\infty\); includes a spurious imaginary branch contribution & \(\infty\) \\
\(\int_{-5}^{0} 1/|x|\,dx\) & not equal to \(\infty\); includes a spurious imaginary branch contribution & \(\infty\) \\
\(\int_{-6}^{0} 1/|x|\,dx\) & not equal to \(\infty\); includes a spurious imaginary branch contribution & \(\infty\) \\
\bottomrule
\end{tabular}
\end{center}

\subsection{Related Issues Assessment}

The bug-finding harness performed a deduplication check.  The closest related
family is improper real integration evaluated through logarithmic
antiderivatives across a branch cut, and this failure mode is already
publicly reported in the open issue
\href{https://github.com/sympy/sympy/issues/23337}{SymPy~\#23337}
(``Spurious I*pi in improper definite integral'').  That issue documents the
same root cause: for \(\texttt{integrate(1/(x-3), (x, 1, 3))}\), the
logarithmic antiderivative \(\log(x-3)\) evaluates to \(\log(2)+i\pi\) at the
finite endpoint via the principal complex log, and this \(i\pi\) branch
constant fails to cancel during definite interval evaluation --- exactly the
\texttt{C = self.subs(x, c)} mechanism diagnosed here.  The deduplication
verdict is therefore a specific new instance of an already-reported defect:
the concrete \(\texttt{integrate(1/Abs(x), (x, -1, 0))}\) reproducer returning
\(\infty+i\pi\) does not appear verbatim in \#23337, but it shares the same
root cause and should be cross-referenced against \#23337 rather than filed as
fully novel.

\subsection{Proposed SymPy Regression Test}

\medskip

\begin{lstlisting}[style=bugpython]
import pytest
from sympy import Abs, integrate, oo, symbols


@pytest.mark.parametrize("a", [1, 2, 3, 4, 5, 6])
def test_improper_integral_abs_negative_side_has_no_imaginary_part(a):
    x = symbols("x", real=True)
    assert integrate(1/Abs(x), (x, -a, 0)) == oo
\end{lstlisting}

\end{document}
